\documentclass{article}

\usepackage[swedish]{babel}
\usepackage{amsmath}
\usepackage{amssymb}

\title{Lösningsformeln för en homogen linjär differentialekvation av ordning 2 med konstanta koefficienter}
\author{Marcus Persson}
\date{\today}

\begin{document}
\maketitle

\section*{\underbar{Sats}}
\noindent Lösningen till en homogen linjär differentialekvation av ordning 2 med konstant koefficienter $a,b \in \mathbb{R}$ till $y''+ay'+by = 0$, där $r_1,r_2$ är röterna till ekvationen, sådan att lösningarna blir:

\begin{equation}
    y = Ae^{r_1x} + Be^{r_2x},
    \quad r_1 \neq r_2,\quad A,B \in \mathbb{R}
\end{equation}

\begin{equation}
    y = (Ax+B)e^{r_1x},
    \quad r_1 = r_2, \quad A,B \in \mathbb{R}
\end{equation}

\section*{\underbar{Bevis}}

Låt $y,z$ vara två deriverbar funktion med avsende på $x$, på något intervall $I$, sådan att:

\begin{equation}
    y = ze^{wx}, 
    \quad w \in \mathbb{R}
\end{equation}

\begin{equation}
\left\{
\begin{aligned}
    y &= ze^{wx},\\
    y' &= (z'+wz)e^{wx},\\
    y'' &= (z''+2wz'+w^2z)e^{wx}.
\end{aligned}
\right.
\end{equation}

\begin{equation}
\begin{gathered}
    y''+ay'+by = 0\\
    \Leftrightarrow\\
    (z''+(2w+a)z'+(w^2+aw+b)z)e^{wx} = 0
\end{gathered}
\end{equation}

\noindent Notera att $e^{wx} \neq 0$, då ser vi till att välja $w=r_1$ sådan att $w^2+aw+b=0$. Vi kan då skriva det som $(w-r_1)(w-r_2)=0$, då följer att $a=-(r_1+r_2)$ och ekvationen blir:

\begin{equation}
\begin{gathered}
    (z''+(2w+a)z'+(w^2+aw+b)z)e^{wx} = 0\\
    \Leftrightarrow\\
    z''+(r_1-r_2)z'=0
\end{gathered}
\end{equation}

\noindent Vi löser det som en homogen första ordnings differentialekvation med konstant koefficienter, där $P(x)=r_1-r_2$ och $Q(x)=0$, sådan att:

\begin{equation}
    \mu(x) = e^{\int P(x)dx} = e^{(r_1-r_2)x} 
\end{equation}

\begin{equation}
\begin{gathered}
    z' = (\mu(x))^{-1}\int \mu(x)Q(x) + (\mu(x))^{-1}C,
    \quad C \in \mathbb{R}\\
    \Leftrightarrow\\
    z' = Ce^{(r_2-r_1)x}\\
\end{gathered}
\end{equation}

\subsection*{\underbar{Fall 1 ($r_1 \neq r_2$):}}

\begin{equation}
    z' = Ce^{(r_2-r_1)x}
\end{equation}

\begin{equation}
\begin{gathered}
    z = \int zdx = \int Ce^{(r_2-r_1)x}dx
    = 
    \frac{C}{r_2-r_1}e^{(r_2-r_1)x} + A,
    \quad A \in \mathbb{R}
\end{gathered}
\end{equation}

\begin{equation}
\begin{gathered}
    y = ze^{r_1x}\\
    \Leftrightarrow\\
    y = Be^{(r_2-r_1)x}e^{r_1x} + Ae^{r_1x},\quad B = \frac{C}{r_2-r_1}\\
    \Leftrightarrow\\
    y = Ae^{r_1x} + Be^{r_2x}
\end{gathered}
\end{equation}

\subsection*{\underbar{Fall 2 ($r_1 = r_2$):}}

\begin{equation}
    z' = A, \quad A = C
\end{equation}

\begin{equation}
    z = \int zdx = \int Adx = Ax + B
\end{equation}

\begin{equation}
\begin{gathered}
    y = ze^{r_1x}\\
    \Leftrightarrow\\
    y = (Ax + B)e^{r_1x}
\end{gathered}
\end{equation}

\noindent Där av $\exists$ lösningsformeln för en homogen linjär differentialekvation av ordning 2 med konstanta koefficienter

\begin{flushright}
    $\blacksquare$
\end{flushright}

\end{document}